% Hiệu
\section{Yêu cầu Hệ thống}
Dựa trên kết quả khảo sát nghiệp vụ ở các mục trước, phần này tổng hợp và phân loại các yêu cầu hệ thống cho nền tảng AR-IMMS thành hai nhóm chính: \textbf{Yêu cầu chức năng} (Functional Requirements) và \textbf{Yêu cầu phi chức năng} (Non-Functional Requirements). Các yêu cầu này là cơ sở trực tiếp để xác định các thực thể dữ liệu, quan hệ và ràng buộc trong thiết kế Cơ sở Dữ liệu ở các chương tiếp theo.

%% =============================================
%% 1.5.1 YÊU CẦU CHỨC NĂNG
%% =============================================
\subsection{Yêu cầu chức năng}
Hệ thống AR-IMMS phục vụ bốn nhóm tác nhân chính: \textit{Quản trị viên Hệ thống} (System Administrator), \textit{Vận hành viên} (System Operator), \textit{Kỹ thuật viên} (Technician) và \textit{Hệ thống Giám sát \& Cảnh báo Tự động} (Automated Monitoring \& Alerting System). Dưới đây là danh sách yêu cầu chức năng phân theo từng vai trò.

% ---- 1.5.1.1 Quản trị viên Hệ thống ----
\subsubsection{Quản trị viên Hệ thống (System Administrator)}
\begin{enumerate}[label=\textbf{FA-\arabic*}., leftmargin=2.5em]
    \item \textbf{Quản lý tài khoản người dùng:} Tạo mới, cập nhật thông tin, vô hiệu hóa và gán vai trò (Administrator, Operator, Technician) cho các tài khoản người dùng trên hệ thống.
    \item \textbf{Cấu hình thiết lập hệ thống:} Thiết lập và chỉnh sửa các thông số cấu hình toàn hệ thống (system-wide settings) bao gồm chu kỳ thu thập dữ liệu, ngưỡng cảnh báo mặc định, chính sách bảo mật và các tham số vận hành khác.
    \item \textbf{Giám sát tổng quan sức khỏe hệ thống:} Xem các chỉ số tổng hợp (aggregated indicators) về tình trạng hoạt động của toàn bộ hệ thống, bao gồm: tính khả dụng của dịch vụ (service availability), trạng thái kết nối của bộ thu thập dữ liệu (data collector connectivity) và lưu lượng cảnh báo (alert volume) trên tất cả các thành phần tích hợp.
    \item \textbf{Quản lý môi trường Mini Data Center:} Quản lý cấu hình các môi trường Mini Data Center thử nghiệm, bao gồm đăng ký site, phòng máy, và các thông số hệ thống liên quan.
    \item \textbf{Truy vấn nhật ký kiểm toán:} Xem và tra cứu nhật ký kiểm toán (audit log) ghi lại toàn bộ lịch sử thay đổi cấu hình, các hành động quản trị, và các sự kiện quan trọng trên hệ thống.
\end{enumerate}

% ---- 1.5.1.2 Vận hành viên ----
\subsubsection{Vận hành viên (System Operator)}
\begin{enumerate}[label=\textbf{FO-\arabic*}., leftmargin=2.5em]
    \item \textbf{Quản lý tài sản hạ tầng:} Đăng ký, cập nhật và quản lý thông tin các tủ Rack, máy chủ (Server/Node) và các khối tải công việc (Workload/Container) tương ứng theo hệ thống phân cấp: Site $\rightarrow$ Room $\rightarrow$ Rack $\rightarrow$ Node $\rightarrow$ Workload.
    \item \textbf{Liên kết mã đánh dấu không gian (Spatial Marker):} Gán mã định danh QR/ArUco Marker cho một máy chủ hoặc node đã đăng ký trong quá trình tiếp nhận thiết bị (onboarding). Cho phép cập nhật liên kết khi phần cứng được di chuyển hoặc thay thế.
    \item \textbf{Giám sát trạng thái vận hành:} Theo dõi trạng thái vận hành theo thời gian thực (real-time) và dữ liệu lịch sử (historical) của toàn bộ hạ tầng Mini Data Center, bao gồm mức sử dụng CPU, RAM, ổ đĩa, lưu lượng mạng và trạng thái container.
    \item \textbf{Tương tác Digital Twin:} Xem và tương tác với mô hình Digital Twin để nhận diện trạng thái, vị trí và mối quan hệ phân cấp giữa các tủ Rack, máy chủ và dịch vụ đang chạy.
    \item \textbf{Cấu hình ngưỡng cảnh báo:} Thiết lập các ngưỡng (threshold) cho từng chỉ số giám sát, rà soát các điều kiện bất thường và xác nhận (acknowledge) các cảnh báo đang hoạt động.
    \item \textbf{Quản lý sự cố và phiếu bảo trì:} Xác nhận cảnh báo đang hoạt động; khi cần xử lý, tạo phiếu sự cố (incident) hoặc phiếu bảo trì (maintenance ticket) tham chiếu đến cảnh báo gốc. Gán, theo dõi phiếu công việc, lập báo cáo vận hành và quản lý lịch sử dịch vụ.
    \item \textbf{Phân tích xu hướng và báo cáo:} Xem các báo cáo xu hướng lịch sử (historical trend reports), các chỉ số quy hoạch dung lượng (capacity planning metrics) và phân tích PUE/tiêu thụ điện năng, được nhóm theo site, phòng máy, tủ rack và node.
    \item \textbf{Quản lý thông số kỹ thuật thiết bị:} Ghi nhận và cập nhật thông số kỹ thuật phần cứng (hardware specifications), thông tin bảo hành (warranty) và lịch sử bảo trì (maintenance history) cho các tài sản đã đăng ký.
\end{enumerate}

% ---- 1.5.1.3 Kỹ thuật viên ----
\subsubsection{Kỹ thuật viên (Technician)}
\begin{enumerate}[label=\textbf{FT-\arabic*}., leftmargin=2.5em]
    \item \textbf{Tiếp nhận phiếu công việc:} Xem danh sách các phiếu sự cố hoặc phiếu bảo trì được giao, nhận diện tủ Rack và máy chủ cần kiểm tra.
    \item \textbf{Nhận diện thiết bị bằng AR:} Sử dụng ứng dụng AR trên thiết bị di động để quét mã QR/ArUco Marker gắn trên phần cứng, từ đó nhận diện thiết bị và truy cập thông tin vận hành theo ngữ cảnh (contextual operational information) được hiển thị trực tiếp trên hình ảnh camera.
    \item \textbf{Xem thông tin chi tiết thiết bị:} Truy cập trạng thái máy chủ, dữ liệu giám sát thời gian thực (telemetry), cảnh báo đang hoạt động, thông tin chẩn đoán (diagnostic information) và danh sách các workload/dịch vụ liên quan.
    \item \textbf{Cập nhật tiến trình xử lý:} Xác nhận tiếp nhận công việc, cập nhật trạng thái xử lý, ghi nhận ghi chú điều tra (investigation notes), nguyên nhân xác định (identified causes) và các biện pháp khắc phục (corrective actions).
    \item \textbf{Hoàn tất và đóng phiếu:} Ghi nhận kết quả bảo trì, xem lại lịch sử dịch vụ của thiết bị, đánh dấu sự cố đã giải quyết (resolved) và gửi yêu cầu đóng phiếu (ticket closure request) để Vận hành viên phê duyệt.
\end{enumerate}

% ---- 1.5.1.4 Hệ thống Giám sát & Cảnh báo Tự động ----
\subsubsection{Hệ thống Giám sát và Cảnh báo Tự động}
\begin{enumerate}[label=\textbf{FS-\arabic*}., leftmargin=2.5em]
    \item \textbf{Thu thập dữ liệu vận hành:} Tự động thu thập dữ liệu telemetry (CPU, RAM, ổ đĩa, mạng, Docker container) và trạng thái kết nối từ các máy chủ đã đăng ký theo chu kỳ được cấu hình.
    \item \textbf{Phát hiện bất thường:} Tự động phát hiện các máy chủ không khả dụng (unavailable), các luồng thu thập dữ liệu bị gián đoạn (interrupted data collection) và các chỉ số vận hành vượt ngưỡng đã cấu hình.
    \item \textbf{Quản lý vòng đời cảnh báo:} Tạo mới và cập nhật trạng thái cảnh báo xuyên suốt vòng đời: Mở (Open) $\rightarrow$ Đã xác nhận (Acknowledged) $\rightarrow$ Đã giải quyết (Resolved) $\rightarrow$ Đã đóng (Closed).
    \item \textbf{Lưu trữ dữ liệu lịch sử và kiểm toán:} Lưu trữ dữ liệu telemetry lịch sử vào cơ sở dữ liệu chuỗi thời gian (time-series database); ghi nhận các hoạt động giám sát, cảnh báo và bảo trì quan trọng vào nhật ký kiểm toán (audit log).
\end{enumerate}

% ---- Bảng tổng hợp yêu cầu chức năng ----
\vspace{0.5em}
\noindent\textbf{Bảng tổng hợp yêu cầu chức năng theo vai trò:}

\vspace{0.3em}
\begin{longtable}{|c|l|p{8.5cm}|}
\hline
\textbf{Mã} & \textbf{Vai trò} & \textbf{Yêu cầu chức năng} \\
\hline
\endfirsthead
\hline
\textbf{Mã} & \textbf{Vai trò} & \textbf{Yêu cầu chức năng} \\
\hline
\endhead
\hline
\endfoot
FA-1 & Administrator & Quản lý tài khoản người dùng (CRUD, gán vai trò) \\
\hline
FA-2 & Administrator & Cấu hình thiết lập toàn hệ thống \\
\hline
FA-3 & Administrator & Giám sát tổng quan sức khỏe hệ thống \\
\hline
FA-4 & Administrator & Quản lý môi trường Mini Data Center \\
\hline
FA-5 & Administrator & Truy vấn nhật ký kiểm toán \\
\hline
FO-1 & Operator & Quản lý tài sản hạ tầng (Rack, Node, Workload) \\
\hline
FO-2 & Operator & Liên kết mã đánh dấu không gian QR/ArUco \\
\hline
FO-3 & Operator & Giám sát trạng thái vận hành thời gian thực \& lịch sử \\
\hline
FO-4 & Operator & Tương tác Digital Twin \\
\hline
FO-5 & Operator & Cấu hình ngưỡng cảnh báo \\
\hline
FO-6 & Operator & Quản lý sự cố và phiếu bảo trì \\
\hline
FO-7 & Operator & Phân tích xu hướng và báo cáo vận hành \\
\hline
FO-8 & Operator & Quản lý thông số kỹ thuật \& bảo hành thiết bị \\
\hline
FT-1 & Technician & Tiếp nhận phiếu công việc được giao \\
\hline
FT-2 & Technician & Nhận diện thiết bị bằng AR (QR/ArUco scan) \\
\hline
FT-3 & Technician & Xem thông tin chi tiết thiết bị \& telemetry \\
\hline
FT-4 & Technician & Cập nhật tiến trình xử lý \\
\hline
FT-5 & Technician & Hoàn tất và đóng phiếu bảo trì \\
\hline
FS-1 & Hệ thống tự động & Thu thập dữ liệu vận hành theo chu kỳ \\
\hline
FS-2 & Hệ thống tự động & Phát hiện bất thường \& vượt ngưỡng \\
\hline
FS-3 & Hệ thống tự động & Quản lý vòng đời cảnh báo \\
\hline
FS-4 & Hệ thống tự động & Lưu trữ telemetry lịch sử \& nhật ký kiểm toán \\
\hline
\end{longtable}

%% =============================================
%% 1.5.2 YÊU CẦU PHI CHỨC NĂNG
%% =============================================
\subsection{Yêu cầu phi chức năng}
Bên cạnh các yêu cầu chức năng, hệ thống AR-IMMS cần đáp ứng các yêu cầu phi chức năng sau đây để đảm bảo chất lượng vận hành trong thực tiễn. Các yêu cầu phi chức năng này cũng ảnh hưởng trực tiếp đến việc thiết kế cấu trúc CSDL, lựa chọn kiểu dữ liệu, chiến lược đánh chỉ mục (indexing), và cơ chế phân vùng (partitioning).

% ---- 1.5.2.1 Hiệu năng (Performance) ----
\subsubsection{Hiệu năng (Performance)}
\begin{itemize}[leftmargin=2em]
    \item Dữ liệu telemetry từ các máy chủ phải được thu thập theo chu kỳ cấu hình được, với chu kỳ mặc định là \textbf{5 giây}.
    \item Một máy chủ sẽ được báo cáo là \textit{không khả dụng} (unavailable) khi không nhận được dữ liệu telemetry trong vòng hơn \textbf{90 giây}.
    \item Hệ thống phải hỗ trợ truyền phát dữ liệu theo thời gian thực qua WebSocket với độ trễ thấp, cho phép điều hướng drilldown liền mạch từ tổng quan đến chi tiết container trong tối đa 3 lần nhấp chuột.
\end{itemize}

% ---- 1.5.2.2 Khả dụng (Availability) ----
\subsubsection{Khả dụng (Availability)}
\begin{itemize}[leftmargin=2em]
    \item Nền tảng phải hoạt động ổn định trong các giờ làm việc dự kiến.
    \item Hệ thống phải thông báo cho người dùng khi một máy chủ hoặc bộ thu thập dữ liệu (data collector) trở nên không khả dụng.
    \item Các dịch vụ quan trọng (Backend API, WebSocket Gateway) cần được theo dõi tính khả dụng liên tục.
\end{itemize}

% ---- 1.5.2.3 Bảo mật (Security) ----
\subsubsection{Bảo mật (Security)}
\begin{itemize}[leftmargin=2em]
    \item Hệ thống phải thực thi xác thực (authentication) cho mọi truy cập và phân quyền dựa trên vai trò (Role-Based Access Control -- RBAC) với ba vai trò: Administrator, Operator và Technician.
    \item Sử dụng \textbf{JWT} (JSON Web Token) cho xác thực API và \textbf{mTLS} (mutual TLS) cho các kết nối từ telemetry agent đến backend.
    \item Tất cả các thay đổi cấu hình, cập nhật phiếu bảo trì và các hành động từ xa (remote actions) phải được ghi nhận vào nhật ký kiểm toán (audit log) với đầy đủ thông tin: người thực hiện, thời điểm, hành động, và đối tượng bị tác động.
\end{itemize}

% ---- 1.5.2.4 Độ tin cậy (Reliability) ----
\subsubsection{Độ tin cậy (Reliability)}
\begin{itemize}[leftmargin=2em]
    \item Các telemetry agent phải tự động thử lại (retry) khi truyền dữ liệu thất bại và tự phục hồi (resume) sau khi kết nối lại thành công.
    \item Hệ thống phải triệt tiêu cảnh báo trùng lặp (duplicate suppression) để ngăn chặn tình trạng bão cảnh báo (alert storm).
    \item Các sự cố phải tuân theo vòng đời trạng thái rõ ràng: Open $\rightarrow$ In-Progress $\rightarrow$ Resolved $\rightarrow$ Closed, được quản lý bằng máy trạng thái (state machine) trong CSDL.
\end{itemize}

% ---- 1.5.2.5 Khả năng mở rộng (Scalability) ----
\subsubsection{Khả năng mở rộng (Scalability)}
\begin{itemize}[leftmargin=2em]
    \item Bản PoC (Proof of Concept) phải chứng minh rằng một node xử lý mới được thêm vào sẽ bắt đầu truyền telemetry và xuất hiện trong Digital Twin trong vòng \textbf{5 phút} kể từ khi đăng ký.
    \item Việc thêm node mới chỉ yêu cầu cấu hình thông tin đăng nhập (credentials) và địa chỉ endpoint, \textbf{không} cần sửa đổi mã nguồn dịch vụ lõi (core service code).
    \item Cấu trúc CSDL phải hỗ trợ mở rộng theo chiều ngang (horizontal scalability) -- cho phép thêm nhiều site, room, rack và node mà không ảnh hưởng đến hiệu năng truy vấn.
\end{itemize}

% ---- 1.5.2.6 Khả năng sử dụng (Usability) ----
\subsubsection{Khả năng sử dụng (Usability)}
\begin{itemize}[leftmargin=2em]
    \item Hệ thống cung cấp giao diện dashboard \textbf{desktop-first} bằng tiếng Anh trên nền tảng Web (Next.js) và ứng dụng AR trên Android (React Native).
    \item Giao diện người dùng phải responsive, dễ tiếp cận (accessible), hiển thị rõ ràng trạng thái sức khỏe máy chủ, chi tiết sự cố, các bước khắc phục và hướng dẫn bảo trì.
    \item Hỗ trợ điều hướng phân cấp liền mạch: Site $\rightarrow$ Room $\rightarrow$ Rack $\rightarrow$ Node $\rightarrow$ Workload với giao diện drilldown trực quan.
\end{itemize}

% ---- 1.5.2.7 Bảo trì & An toàn (Maintainability & Safety) ----
\subsubsection{Bảo trì và An toàn (Maintainability \& Safety)}
\begin{itemize}[leftmargin=2em]
    \item Hệ thống phải hỗ trợ kiểm thử tự động (automated testing), ghi log tập trung (centralized logging), và cho phép cấu hình linh hoạt các chỉ số giám sát cùng ngưỡng cảnh báo.
    \item Các hành động có tính phá hủy hoặc ảnh hưởng lớn (ví dụ: khởi động lại dịch vụ từ phiên AR) phải yêu cầu \textbf{xác nhận} (confirmation) trước khi thực thi và được ghi vào nhật ký kiểm toán.
    \item Mã nguồn backend và frontend phải được tổ chức theo kiến trúc module, hỗ trợ triển khai bằng container (containerized deployment) để thuận tiện cho việc bảo trì và nâng cấp.
\end{itemize}

% ---- Bảng tổng hợp yêu cầu phi chức năng ----
\vspace{0.5em}
\noindent\textbf{Bảng tổng hợp yêu cầu phi chức năng:}

\vspace{0.3em}
\begin{longtable}{|c|p{3cm}|p{9cm}|}
\hline
\textbf{Mã} & \textbf{Thuộc tính} & \textbf{Mô tả yêu cầu} \\
\hline
\endfirsthead
\hline
\textbf{Mã} & \textbf{Thuộc tính} & \textbf{Mô tả yêu cầu} \\
\hline
\endhead
\hline
\endfoot
NF-1 & Hiệu năng & Thu thập telemetry mặc định 5s; phát hiện unavailable sau 90s \\
\hline
NF-2 & Hiệu năng & Truyền phát real-time qua WebSocket, drilldown $\leq$ 3 click \\
\hline
NF-3 & Khả dụng & Hoạt động ổn định giờ làm việc; thông báo khi server/collector unavailable \\
\hline
NF-4 & Bảo mật & Xác thực JWT, phân quyền RBAC (3 vai trò), mTLS cho agent \\
\hline
NF-5 & Bảo mật & Audit log đầy đủ cho mọi thay đổi cấu hình và hành động quan trọng \\
\hline
NF-6 & Độ tin cậy & Agent tự retry/resume; triệt tiêu cảnh báo trùng lặp \\
\hline
NF-7 & Độ tin cậy & Vòng đời sự cố: Open $\rightarrow$ In-Progress $\rightarrow$ Resolved $\rightarrow$ Closed \\
\hline
NF-8 & Khả năng mở rộng & Node mới streaming telemetry $\leq$ 5 phút; không sửa mã nguồn lõi \\
\hline
NF-9 & Khả năng sử dụng & Dashboard desktop-first (Web), AR app (Android); giao diện responsive \\
\hline
NF-10 & Bảo trì \& An toàn & Automated testing, centralized logging; xác nhận trước hành động phá hủy \\
\hline
\end{longtable}

%% =============================================
%% 1.5.3 ÁNH XẠ YÊU CẦU VÀO ĐỐI TƯỢNG DỮ LIỆU
%% =============================================
\subsection{Ánh xạ Yêu cầu vào Đối tượng Dữ liệu chính}
Từ các yêu cầu chức năng và phi chức năng trên, có thể nhận diện các nhóm đối tượng dữ liệu chính mà CSDL của hệ thống AR-IMMS cần quản lý. Bảng dưới đây tóm tắt ánh xạ từ yêu cầu sang các thực thể dữ liệu trọng yếu, làm cơ sở cho việc xây dựng mô hình ERD ở Chương 2.

\vspace{0.3em}
\begin{longtable}{|p{3.5cm}|p{4cm}|p{6cm}|}
\hline
\textbf{Nhóm yêu cầu} & \textbf{Đối tượng dữ liệu} & \textbf{Mô tả} \\
\hline
\endfirsthead
\hline
\textbf{Nhóm yêu cầu} & \textbf{Đối tượng dữ liệu} & \textbf{Mô tả} \\
\hline
\endhead
\hline
\endfoot
Quản lý tài khoản, phân quyền (FA-1, NF-4) & User, Role & Tài khoản người dùng, vai trò RBAC, token xác thực \\
\hline
Quản lý tài sản hạ tầng (FO-1, FO-2, FO-8) & Site, Room, Rack, Node, Workload, SpatialMarker & Hệ thống phân cấp tài sản, thông số kỹ thuật, mã marker \\
\hline
Thu thập \& giám sát (FO-3, FS-1, NF-1) & TelemetryRecord, CollectorStatus & Dữ liệu telemetry chuỗi thời gian, trạng thái kết nối agent \\
\hline
Cảnh báo (FO-5, FS-2, FS-3, NF-6, NF-7) & Alert, AlertThreshold & Ngưỡng cảnh báo, cảnh báo với vòng đời trạng thái \\
\hline
Sự cố \& phiếu bảo trì (FO-6, FT-1 -- FT-5) & Incident, Ticket, TicketActivity & Phiếu sự cố, phiếu bảo trì, ghi chú xử lý, lịch sử dịch vụ \\
\hline
Bảo hành \& lịch sử bảo trì (FO-8) & Warranty, MaintenanceLog & Thông tin bảo hành, lịch sử bảo trì thiết bị \\
\hline
Kiểm toán (FA-5, NF-5, FS-4) & AuditLog & Nhật ký kiểm toán ghi lại mọi thay đổi quan trọng \\
\hline
Báo cáo \& phân tích (FO-7, NF-2) & Report, TrendAnalytics & Báo cáo vận hành, phân tích xu hướng, PUE metrics \\
\hline
Cấu hình hệ thống (FA-2, FA-4) & SystemConfig, DataCenterConfig & Tham số cấu hình hệ thống, thiết lập môi trường DC \\
\hline
\end{longtable}

\vspace{0.3em}
\noindent Các đối tượng dữ liệu được xác định ở trên sẽ được phân tích chi tiết, xác định thuộc tính, mối quan hệ và ràng buộc toàn vẹn trong \textbf{Chương 2 -- Thiết kế Mô hình Dữ liệu quan niệm} và tiếp tục được chuẩn hóa trong \textbf{Chương 3 -- Thiết kế Mô hình Logic \& Chuẩn hóa}.

